\documentclass[12pt]{article}
\usepackage[margin=0.75in]{geometry}
\usepackage{fontspec}
\setmainfont{Latin Modern Roman}
\usepackage{microtype}
\usepackage{fancyhdr}
\usepackage{titlesec}
\usepackage{enumitem}
\usepackage{xcolor}
\usepackage[hidelinks]{hyperref}
\setlength{\parindent}{1.1em}
\setlength{\parskip}{0.35em}
\setlength{\headheight}{15pt}
\titleformat{\section}{\large\bfseries\scshape}{}{0pt}{}
\titleformat{\subsection}{\normalsize\bfseries}{}{0pt}{}
\pagestyle{fancy}
\fancyhf{}
\fancyhead[L]{Whately Elements of Logic}
\fancyhead[R]{Teacher Guide}
\fancyfoot[C]{\thepage}
\renewcommand{\headrulewidth}{0.4pt}

\begin{document}
\begin{center}
{\LARGE\bfseries Teacher Guide: Week 5}\par\vspace{0.25em}
{\large\scshape The Shape of an Argument}\par\vspace{0.4em}
{Mood, Figure, and Formal Validity in Whately and \textit{Julius Caesar}}\par\vspace{0.6em}
\rule{0.82\textwidth}{0.4pt}
\end{center}

\section*{Week 5 Overview}
Week 5 is the second syllogism week. Week 4 taught students to map premises, conclusion, and major/minor/middle terms in ordinary language. Week 5 adds a light formal layer: mood, figure, and formal validity. Students classify simple categorical syllogisms by A/E/I/O pattern and by the placement of the middle term. The aim is not full memorization of every traditional valid mood. The aim is disciplined judgment: a student should be able to say whether the shape of an argument carries its conclusion.

The lit-anchor stays with \textit{Julius Caesar}. Brutus's funeral speech gives a clean-looking argument whose form can be mapped. Antony's speech then presses on the middle term, ``ambitious,'' and asks whether Caesar truly belongs under it. This prepares Week 6, where students will distinguish validity from truth and soundness more explicitly.

\section*{Essential Question}
Can the shape of an argument tell us whether the conclusion follows?

\section*{Student-Facing Materials}
\begin{itemize}[leftmargin=*]
\item \textit{The Shape of an Argument} -- student reading handout.
\item \textit{Week 5 Argument Lab: Mood, Figure, and Formal Validity} -- student lab for classifying simple syllogisms and producing a portfolio classification card.
\item \textit{Julius Caesar Lit-Anchor: Brutus, Antony, and the Shape of Proof} -- Shakespeare anchor passage for testing the form of public arguments about Caesar's ambition.
\end{itemize}

\section*{Learning Goals}
Students will be able to:
\begin{enumerate}[leftmargin=*]
\item Distinguish a form test from a fact test.
\item Identify A, E, I, and O propositions.
\item Define mood as the A/E/I/O order of major premise, minor premise, and conclusion.
\item Define figure as the placement of the middle term in the premises.
\item Identify S, P, and M in a categorical syllogism.
\item Recognize the valid beginner form A-A-A in the first figure: All M are P; All S are M; therefore All S are P.
\item Explain the problem of the undistributed middle in ordinary language.
\item Explain why two negative premises do not connect S and P.
\item Apply mood-and-figure testing to Brutus and Antony in \textit{Julius Caesar}.
\end{enumerate}

\section*{Key Vocabulary}
mood, figure, formal validity, A proposition, E proposition, I proposition, O proposition, distribution, undistributed middle, illicit process, affirmative conclusion, negative conclusion, form test, fact test

\section*{Core Teaching Emphasis}
The most important distinction is this: \textit{valid} does not mean \textit{true}. A valid form shows that the conclusion follows if the premises are true. It does not certify the premises. Students should not leave Week 5 thinking that a formal map proves the whole argument. They should leave knowing how to separate two questions:
\begin{enumerate}[leftmargin=*]
\item Does the conclusion follow from these premises?
\item Are these premises true or adequately supported?
\end{enumerate}

Use Brutus as the central literary example. His argument can be given a valid A-A-A shape, but the premise ``Caesar's ambition would enslave Rome'' still needs proof. Antony's speech is powerful because it does not merely emote; it contests the middle term that Brutus needs.

\section*{Weekly Lesson Sequence}

\subsection*{Day 1: From Practical Maps to Form}
\begin{itemize}[leftmargin=*]
\item Review Week 4: conclusion, major term, minor term, middle term, major premise, minor premise.
\item Introduce the form test vs. the fact test.
\item Work through the whale or careful-proofs example from the reading.
\item Have students mark S, P, and M before discussing mood.
\end{itemize}

\subsection*{Day 2: Mood and A/E/I/O Review}
\begin{itemize}[leftmargin=*]
\item Review A/E/I/O proposition forms from Week 3.
\item Show how a syllogism's mood is formed by listing the proposition types in order.
\item Practice with three quick examples: A-A-A, E-A-E, and an invalid A-A-A with the middle in the wrong position.
\item Keep the emphasis on explanation, not naming every traditional mnemonic.
\end{itemize}

\subsection*{Day 3: Figure and Invalid Forms}
\begin{itemize}[leftmargin=*]
\item Introduce figure as placement of the middle term.
\item Diagram All M are P; All S are M; therefore All S are P.
\item Contrast it with All P are M; All S are M; therefore All S are P.
\item Use the ``careful readers / suspicious readers'' example to explain undistributed middle.
\item Students complete Part III of the lab and repair one invalid form.
\end{itemize}

\subsection*{Day 4: Julius Caesar Lit-Anchor}
\begin{itemize}[leftmargin=*]
\item Read Brutus's Act 3, Scene 2 speech aloud.
\item Build the guided map on the board: All men whose ambition would enslave Rome deserve death; Caesar was such a man; therefore Caesar deserved death.
\item Ask students to distinguish form strength from factual support.
\item Read Antony's answer. Mark how Antony attacks the middle term ``ambitious.''
\item Students begin the \textit{Julius Caesar Mood-and-Figure Card} portfolio artifact.
\end{itemize}

\subsection*{Day 5: Portfolio Card and Synthesis}
\begin{itemize}[leftmargin=*]
\item Students finish either the Week 5 lab card or the lit-anchor card.
\item Short seminar: Which speaker gives the cleaner formal argument? Which speaker gives stronger evidence?
\item Board synthesis: Brutus may have form; Antony attacks the premise; Week 6 will ask how validity differs from soundness.
\item Exit writing: ``A valid form is not enough because...'' Require one example from Brutus, Antony, or the lab.
\end{itemize}

\section*{Answer Guidance: Brutus}
A strong formal reconstruction:
\begin{itemize}[leftmargin=*]
\item \textbf{Major premise:} All men whose ambition would enslave Rome deserve death for Rome's freedom.
\item \textbf{Minor premise:} Caesar was a man whose ambition would enslave Rome.
\item \textbf{Conclusion:} Caesar deserved death for Rome's freedom.
\item \textbf{Mood:} A-A-A.
\item \textbf{Figure:} All M are P; All S are M; therefore All S are P.
\item \textbf{Form judgment:} Valid, if the premises are granted.
\item \textbf{Premise pressure point:} the minor premise. Brutus must prove Caesar belongs under the middle term ``men whose ambition would enslave Rome.''
\end{itemize}

Students may phrase the major term differently, such as ``people deserving death for Rome's freedom.'' The core point is that the conclusion predicate becomes P, Caesar becomes S, and the ambition/slavery class becomes M.

\section*{Answer Guidance: Antony}
Antony's reasoning is partly inductive and rhetorical, so do not force every line into one tidy syllogism. Still, students can map one implied argument:
\begin{itemize}[leftmargin=*]
\item \textbf{Major premise:} No man who enriches Rome, weeps for the poor, and refuses a crown is clearly ambitious in the tyrannical sense.
\item \textbf{Minor premise:} Caesar enriched Rome, wept for the poor, and refused a crown.
\item \textbf{Conclusion:} Caesar is not clearly ambitious in the tyrannical sense.
\item \textbf{Possible mood:} E-A-E if phrased as ``No M are P; Caesar is M; therefore Caesar is not P.''
\item \textbf{Form judgment:} The form can be made valid, but Antony's larger speech also relies on examples, irony, emotion, and crowd movement.
\end{itemize}

The important teaching point is that Antony pressures Brutus's middle term. He does not need to prove Caesar was perfect; he needs to make the crowd doubt that Caesar belongs under the term ``ambitious tyrant.''

\section*{Common Student Errors}
\begin{itemize}[leftmargin=*]
\item Saying ``valid'' when they mean ``persuasive'' or ``morally convincing.''
\item Saying ``invalid'' when they only mean ``one premise is false or doubtful.''
\item Naming mood before identifying the conclusion.
\item Losing track of S and P because the sentence order changes.
\item Treating Antony's irony as if it were a straightforward endorsement of Brutus.
\item Forgetting that Antony's evidence attacks the minor premise in Brutus's argument.
\end{itemize}

\section*{Connection to Future Weeks}
Week 6 should use the Brutus/Antony contrast again. Week 5 can say: Brutus's argument may be formally valid if reconstructed charitably. Week 6 will say: even a valid form does not make an argument sound unless its premises are true. Antony's speech prepares that distinction by making the crowd doubt the premise that Caesar was tyrannically ambitious.

\end{document}
